\documentclass[11pt,a4paper]{article}
\usepackage{kotex}
\usepackage[margin=25mm]{geometry}
\usepackage{hyperref}
\usepackage{xcolor}
\usepackage[most]{tcolorbox} % 'most' for more tcolorbox features
\usepackage{booktabs}
\usepackage{graphicx}
\usepackage{amsmath}
\usepackage{amssymb}
\usepackage{listings}
\usepackage{setspace}
\usepackage{fancyhdr}
\usepackage{adjustbox} % For resizing tables

% PDF Metadata
\hypersetup{
  pdftitle={CSCI103 Lecture 12: 데이터 및 AI 거버넌스},
  pdfauthor={UIUC Student},
  pdfsubject={CSCI E-103: Introduction to the Intellectual Enterprises of CS}
}

% Line spacing
\onehalfspacing % or \linespread{1.25}

% Paragraph settings
\setlength{\parskip}{0.6em}
\setlength{\parindent}{0pt}

% Table row stretch
\renewcommand{\arraystretch}{1.2}

% Listings settings for code blocks
\lstset{
  basicstyle=\ttfamily\small,
  breaklines=true,
  numbers=left,
  numbersep=6pt,
  frame=single,
  captionpos=b, % Caption at the bottom
  keywordstyle=\color{blue},
  stringstyle=\color{red},
  commentstyle=\color{green!50!black},
  showstringspaces=false
}

% Custom tcolorbox styles
\tcbset{
    mybox/.style={
        colback=blue!5!white,
        colframe=blue!75!black,
        fonttitle=\bfseries,
        boxsep=5pt,
        arc=4pt,
        outer arc=4pt,
        left=6pt, right=6pt, top=6pt, bottom=6pt
    },
    summarybox/.style={
        mybox,
        colback=green!5!white,
        colframe=green!75!black,
        title={요약}
    },
    cautionbox/.style={
        mybox,
        colback=red!5!white,
        colframe=red!75!black,
        title={주의}
    },
    examplebox/.style={
        mybox,
        colback=yellow!5!white,
        colframe=yellow!75!black,
        title={예시}
    },
    keyconceptbox/.style={
        mybox,
        colback=purple!5!white,
        colframe=purple!75!black,
        title={핵심 개념}
    }
}

% Header and Footer
\pagestyle{fancy}
\fancyhf{} % Clear all header and footer fields
\fancyhead[L]{CSCI103 Lecture 12: 데이터 및 AI 거버넌스}
\fancyhead[R]{\today}
\fancyfoot[C]{\thepage}

\begin{document}

\begin{tcolorbox}[summarybox]
\textbf{개요} \\
이 문서는 \texttt{CSCI103 Lecture 12} 강의 내용을 바탕으로 데이터 및 AI 거버넌스의 핵심 개념을 다룹니다. \\
정보 거버넌스의 광범위한 이론적 배경부터 데이터 및 AI 거버넌스의 구체적인 원칙, \\
그리고 Databricks의 Unity Catalog를 활용한 실제 구현 방안까지 상세히 설명합니다. \\
거버넌스 성숙도 모델, 위험 관리, 역할 및 책임, 그리고 Databricks 플랫폼의 주요 기능을 통해 \\
조직이 데이터를 안전하고 효율적으로 관리하며 AI를 윤리적으로 활용하는 방법을 이해하는 것이 목표입니다. \\
이 문서는 비전공자도 완벽히 이해할 수 있도록 쉬운 설명과 풍부한 예시를 제공합니다.
\end{tcolorbox}

\newpage
\section*{용어 정리}

다음 표는 강의에서 다루는 주요 용어들을 비전공자도 이해하기 쉽게 설명합니다.

\begin{adjustbox}{max width=\textwidth}
\begin{tabular}{llll}
\toprule
\textbf{용어} & \textbf{쉬운 설명} & \textbf{원어} & \textbf{비고} \\
\midrule
정보 거버넌스 & 조직의 모든 정보 자산을 관리하는 큰 틀 & Information Governance & 데이터 및 AI 거버넌스를 포함하는 상위 개념 \\
데이터 거버넌스 & 데이터의 수집, 저장, 보안 등 관리 규칙 & Data Governance & 전통적인 데이터 관리 방식 \\
AI 거버넌스 & AI 시스템의 개발, 배포, 감독 규칙 & AI Governance & 모델의 윤리, 편향, 책임 등 고려 \\
정책 & 조직이 지켜야 할 일반적인 규칙 & Policy & 법률이나 규제에 의해 정의될 수 있음 \\
절차 & 정책을 실행하기 위한 단계별 방법 & Procedure & 특정 작업을 수행하는 구체적인 과정 \\
표준 & 정책 준수를 위한 기술적 요구사항 & Standard & 데이터 암호화, 접근 제어 방식 등 \\
통제 & 정책 준수를 강제하는 메커니즘 & Control & 다단계 인증, 접근 권한 설정 등 \\
Unity Catalog (UC) & Databricks의 통합 데이터 및 AI 거버넌스 솔루션 & Unity Catalog & 모든 데이터 자산의 메타데이터를 중앙에서 관리 \\
속성 기반 접근 제어 & 데이터의 속성(태그)에 따라 접근 권한을 부여 & Attribute-Based Access Control (ABAC) & 특정 태그가 있는 데이터에만 규칙 적용 \\
데이터 분류 & 데이터 내 민감 정보(PII)를 자동으로 식별 & Data Classification & AI 모델을 활용하여 자동 태그 부여 \\
데이터 리니지 & 데이터의 출처와 흐름을 추적하는 기능 & Data Lineage & 데이터가 어디서 와서 어디로 가는지 파악 \\
레이크하우스 페더레이션 & 외부 데이터베이스를 Databricks에서 통합 관리 & Lakehouse Federation & 다양한 데이터 소스에 대한 통합 거버넌스 제공 \\
델타 쉐어링 & Databricks 자산을 안전하게 공유하는 기술 & Delta Sharing & 데이터, 노트북, AI 모델 등을 외부와 공유 \\
PII & 개인을 식별할 수 있는 정보 & Personally Identifiable Information & 이름, 주소, SSN, 이메일 등 \\
RBAC & 역할에 따라 접근 권한을 부여하는 방식 & Role-Based Access Control & 직무에 필요한 최소한의 권한만 부여 \\
LLM & 대규모 언어 모델 & Large Language Model & 방대한 데이터로 학습된 AI 모델 (예: ChatGPT) \\
\bottomrule
\end{tabular}
\end{adjustbox}
\label{tab:glossary}

\newpage
\section*{핵심 개념 및 원리}

\subsection*{정보 거버넌스 (Information Governance)}

\begin{keyconceptbox}
\textbf{한 줄 핵심 요약:} 정보 거버넌스는 조직이 보유한 모든 정보 자산을 효과적으로 관리하고 보호하기 위한 포괄적인 틀입니다.
\end{keyconceptbox}

\textbf{직관적 예시/비유:}
정보 거버넌스는 마치 한 도시의 모든 건물(데이터)과 도로(정보 흐름)를 관리하는 도시 계획과 같습니다. 어떤 건물을 어디에 지을지, 누가 접근할 수 있는지, 어떤 규제를 따라야 하는지 등 도시 전체의 질서와 안전을 책임지는 큰 그림을 그리는 것입니다. 데이터 거버넌스와 AI 거버넌스는 이 도시 계획 안에서 특정 구역(데이터)이나 특정 유형의 시설(AI 시스템)을 관리하는 더 세부적인 계획에 해당합니다.

\textbf{기술적 정확한 설명:}
정보 거버넌스는 조직 내에서 생성, 수집, 저장, 처리, 공유되는 모든 데이터, 지식, AI 자산을 관리하기 위한 포괄적인 용어입니다. 이는 클라우드 상의 디지털 정보뿐만 아니라 인쇄물, 물리적 문서, 심지어 노트북 잠금 규칙과 같은 물리적 보안 통제까지 포함하는 매우 광범위한 개념입니다.

\subsubsection*{정보 거버넌스의 목적}
정보 거버넌스의 주요 목적은 다음과 같습니다.
\begin{itemize}
    \item \textbf{자산의 효과적인 관리:} 정보 자산을 효율적으로 활용하고 가치를 극대화합니다.
    \item \textbf{데이터의 보안, 정확성, 접근성 보장:} 민감한 정보가 안전하게 보호되고, 정확하며, 필요한 사람이 적절하게 접근할 수 있도록 합니다.
    \item \textbf{법적 및 규제 요구사항 준수:} HIPAA, PCI, GDPR과 같은 법률 및 규제를 준수하여 조직이 법적 문제나 벌금을 피하도록 합니다. 감사(audit) 시 문제가 발생하지 않도록 대비하는 것이 중요합니다.
\end{itemize}

\subsubsection*{정보 거버넌스의 목표}
정보 거버넌스는 크게 두 가지 상충하는 목표를 동시에 달성하고자 합니다: 정보 가치 극대화와 위험 완화.

\begin{figure}[h!]
    \centering
    \includegraphics[width=0.8\textwidth]{example-quadrant.png} % Placeholder for a conceptual diagram
    \caption{정보 거버넌스의 주요 목표 (개념도)}
    \label{fig:info_gov_goals}
\end{figure}

\textbf{정보 거버넌스의 4가지 주요 목표:}
\begin{enumerate}
    \item \textbf{정보 가치 극대화 (Maximize Information Value):}
    데이터는 고갈되지 않는 재생 가능한 자원이며, 지속적으로 생성됩니다. 조직은 이 데이터를 통해 최대의 가치를 창출하고자 합니다.
    \item \textbf{위험 완화 (Mitigate Risk):}
    동시에 데이터 유출, 오용, 규제 미준수 등으로 인한 위험을 최소화해야 합니다. 가치 극대화가 위험을 초래하지 않도록 균형을 맞추는 것이 중요합니다.
    \item \textbf{정보 보안 강화 (Enhance Information Security):}
    최신 보안 표준(예: 하드웨어 기반 인증 키)을 도입하여 조직의 정보 자산을 보호합니다.
    \item \textbf{정보 수명 주기 관리 최적화 (Optimize Information Life Cycle Management):}
    데이터의 수집부터 저장, 보관, 폐기까지의 전 과정을 효율적으로 관리합니다. 오래된 데이터의 보관 정책, 아카이빙 전략 등을 포함합니다.
    \item \textbf{투명성 및 책임성 증진 (Promote Transparency and Accountability):}
    데이터 관리 및 사용에 대한 명확한 역할과 책임을 설정하고, 감사 추적(audit trail) 및 보고 메커니즘을 구축하여 투명성을 확보합니다.
\end{enumerate}

\subsubsection*{정보 거버넌스 프레임워크}
정보 거버넌스는 조직의 최상위 리더십부터 실제 데이터를 다루는 현업 부서까지 전반에 걸쳐 정의되고 운영됩니다.

\begin{figure}[h!]
    \centering
    \includegraphics[width=0.9\textwidth]{example-framework.png} % Placeholder for a conceptual diagram
    \caption{정보 거버넌스 정의 및 운영 프레임워크 (개념도)}
    \label{fig:info_gov_framework}
\end{figure}

\textbf{1. 정의 (Definition) - 상위 레벨:}
\begin{itemize}
    \item \textbf{경영진 리더십 (Executive Leadership):} CEO, COO, CGO(Chief Governance Officer) 등이 기업 표준 및 정책을 정의합니다.
    \item \textbf{법률 및 규제 준수 부서 (Legal and Compliance):} 기업 거버넌스 사무실을 지원하며, 외부 법률 및 규제(예: GDPR, CCPA)를 해석하고 내부 정책에 반영합니다.
    \item \textbf{사업부 리더십 (Business Unit Leadership):} 기업 정책을 준수하면서 각 사업부의 특성에 맞는 표준 및 정책을 정의합니다.
    \item \textbf{IT 표준 및 정책 (IT Standards and Policies):} 사업 목표를 지원하며, 모든 정책을 준수하는 IT 표준을 수립합니다.
    \item \textbf{데이터/거버넌스 전문 사무실 (Data/Governance Specialty Office):} 실제 데이터 자산(테이블, 파일, 모델 등)에 대한 표준을 정의합니다.
    \item \textbf{외부 지침 (External Guidelines):} 법률, 산업 규제, 인증, 문화적/윤리적 권고 사항 등 외부 요인이 모든 정책 정의에 영향을 미칩니다.
\end{itemize}

\textbf{2. 운영 (Operations) - 하위 레벨:}
\begin{itemize}
    \item \textbf{IT 및 사업부 (IT and Business Units):} 정의된 규칙을 실제 플랫폼과 데이터/AI 자산에 구현합니다.
    \item \textbf{플랫폼 및 데이터/AI 자산 (Platforms and Data/AI Assets):} Databricks와 같은 도구를 사용하여 거버넌스 정책을 기술적으로 적용합니다. Databricks는 기업 표준 및 정책을 정의하는 것이 아니라, IT 조직과 사업부가 규제 준수를 위해 필요한 기능을 구현하도록 돕습니다.
\end{itemize}

\begin{cautionbox}
\textbf{데이터 유출과 거버넌스:}
DoorDash나 Capital One과 같은 데이터 유출 사고는 직접적으로 거버넌스 다이어그램에 명시되지는 않지만, '법률 및 규제 준수'와 '조직에 미치는 영향'이라는 측면에서 중요한 동인으로 작용합니다. 이러한 사고는 금전적 손실뿐만 아니라 기업의 평판 위험(reputational risk)과 소비자 신뢰 상실로 이어지기 때문에, 조직은 이를 방지하기 위해 거버넌스 정책을 강화하게 됩니다.
\end{cautionbox}

\subsubsection*{거버넌스의 4가지 구성 요소: 정책, 절차, 표준, 통제}
거버넌스는 이 네 가지 요소가 유기적으로 결합하여 작동하는 총체적인 접근 방식입니다.

\begin{figure}[h!]
    \centering
    \includegraphics[width=0.7\textwidth]{example-holistic.png} % Placeholder for a conceptual diagram
    \caption{거버넌스의 4가지 구성 요소 (개념도)}
    \label{fig:gov_holistic}
\end{figure}

\textbf{1. 정책 (Policies):}
\begin{itemize}
    \item \textbf{정의:} 조직이 달성하고자 하는 목표와 행동 원칙을 명시한 상위 수준의 규칙입니다.
    \item \textbf{예시:} 데이터 보호 정책, 개인정보 보호 정책, 데이터 보존 정책. (예: GDPR은 개인 데이터 보호에 대한 정책을 제시합니다.)
\end{itemize}

\textbf{2. 절차 (Procedures):}
\begin{itemize}
    \item \textbf{정의:} 정책을 구현하기 위한 단계별 지침 또는 프로세스입니다.
    \item \textbf{예시:} 개인 데이터 접근 요청 처리 절차 (예: 개인이 자신의 데이터 접근 또는 삭제를 요청했을 때 45일 이내에 처리해야 하는 단계별 지침).
\end{itemize}

\textbf{3. 표준 (Standards):}
\begin{itemize}
    \item \textbf{정의:} 정책 및 절차를 준수하기 위한 기술적 또는 운영적 요구사항입니다.
    \item \textbf{예시:} 데이터 암호화 표준, 데이터 최소화 원칙, 역할 기반 접근 제어(RBAC) 구현 표준. (예: 민감 데이터는 특정 방식으로 암호화해야 하며, 사용자에게는 직무에 필요한 최소한의 권한만 부여해야 합니다.)
\end{itemize}

\textbf{4. 통제 (Controls):}
\begin{itemize}
    \item \textbf{정의:} 정책 및 표준이 실제로 준수되도록 강제하는 메커니즘 또는 조치입니다.
    \item \textbf{예시:}
    \begin{itemize}
        \item \textbf{기술적 통제:} 다단계 인증(MFA), 접근 권한 관리 시스템.
        \item \textbf{관리적 통제:} 접근 승인 프로세스, 보안 교육.
        \item \textbf{물리적 통제:} 서버실 출입 통제, 노트북 잠금 규칙.
    \end{itemize}
\end{itemize}

\begin{examplebox}
\textbf{GDPR (일반 데이터 보호 규정) 예시:}
\begin{itemize}
    \item \textbf{정책:} "개인의 데이터는 보호되어야 하며, 삭제 요청 시 즉시 처리되어야 한다." (데이터 보호 정책)
    \item \textbf{절차:} "개인 데이터 삭제 요청 접수 시, 담당 부서는 45일 이내에 해당 데이터를 시스템에서 완전히 삭제하고, 삭제 확인서를 발급한다." (데이터 삭제 요청 처리 절차)
    \item \textbf{표준:} "개인 데이터는 전송 및 저장 시 AES-256 암호화를 적용해야 하며, 접근은 역할 기반 접근 제어(RBAC) 원칙에 따라 최소한의 인원에게만 허용한다." (암호화 및 RBAC 표준)
    \item \textbf{통제:} "데이터베이스 접근 시 다단계 인증(MFA)을 필수로 적용하고, 모든 데이터 삭제 작업은 감사 로그에 기록되어야 한다." (기술적 통제 및 관리적 통제)
\end{itemize}
\end{examplebox}

\newpage
\subsection*{데이터 거버넌스 vs. AI 거버넌스}

\begin{keyconceptbox}
\textbf{한 줄 핵심 요약:} 데이터 거버넌스는 데이터 자체의 관리와 보안에 중점을 두는 반면, AI 거버넌스는 AI 시스템의 개발, 배포, 운영 전반에 걸쳐 윤리적이고 책임감 있는 사용을 보장하는 데 초점을 맞춥니다.
\end{keyconceptbox}

\textbf{직관적 예시/비유:}
\begin{itemize}
    \item \textbf{데이터 거버넌스:} 금고에 돈(데이터)을 안전하게 보관하고, 누가 얼마를 인출할 수 있는지 규칙을 정하는 것과 같습니다. 돈 자체의 보관과 접근 권한에 대한 이야기입니다.
    \item \textbf{AI 거버넌스:} 금고에서 인출한 돈으로 어떤 사업을 할지, 그 사업이 사회에 해를 끼치지는 않을지, 공정하게 운영될지 등을 고민하는 것과 같습니다. 돈을 사용하는 '행위'와 그 '결과'에 대한 이야기입니다. 특히 최근에는 대규모 언어 모델(LLM)처럼 출처가 불분명한 데이터로 학습된 모델을 사용할 때, 그 모델이 생성하는 결과에 대한 책임과 윤리적 문제가 더욱 중요해졌습니다.
\end{itemize}

\textbf{기술적 정확한 설명:}
\begin{itemize}
    \item \textbf{데이터 거버넌스 (Data Governance):} 데이터의 수집, 저장, 보안, 품질, 보존 등 데이터 수명 주기 전반에 걸쳐 정책, 절차, 표준을 수립하고 적용하는 전통적인 접근 방식입니다. 조직의 데이터 자산을 보호하고 가치를 극대화하는 데 중점을 둡니다.
    \item \textbf{AI 거버넌스 (AI Governance):} AI 시스템(모델, 알고리즘)의 개발, 배포, 감독에 대한 정책, 절차, 표준을 다룹니다. AI 모델 학습에 사용된 데이터의 적법성, 모델의 편향성, 투명성, 설명 가능성, 윤리적 사용, 그리고 AI 시스템으로 인한 잠재적 위험을 관리하는 것이 핵심입니다.
\end{itemize}

\subsubsection*{데이터 거버넌스와 AI 거버넌스의 관계}
\begin{cautionbox}
\textbf{AI 거버넌스는 데이터 거버넌스 없이는 불가능합니다.}
AI 모델은 데이터에 기반하여 학습되고 작동하므로, 데이터가 제대로 거버넌스되지 않으면 AI 거버넌스도 효과적으로 이루어질 수 없습니다. 데이터의 보안, 품질, 적법성이 확보되지 않은 상태에서 AI 시스템의 윤리적 사용을 논하는 것은 어렵습니다. 이는 조직의 성숙도 곡선에서 데이터 거버넌스가 AI 거버넌스보다 선행되어야 함을 의미합니다.
\end{cautionbox}

\subsubsection*{데이터 거버넌스 및 AI 거버넌스 역량}
조직이 데이터 및 AI 거버넌스를 구축할 때 고려해야 할 주요 역량은 다음과 같습니다.

\begin{adjustbox}{max width=\textwidth}
\begin{tabular}{p{0.2\textwidth}p{0.35\textwidth}p{0.35\textwidth}}
\toprule
\textbf{영역} & \textbf{데이터 거버넌스 역량} & \textbf{AI 거버넌스 역량} \\
\midrule
\textbf{자산 관리} & 메타데이터 관리, 데이터 검색, 마스터 데이터 관리(MDM), 데이터 분류 & 모델 관리 및 수명 주기, 모델 모니터링, 모델 버전 관리 \\
\textbf{정책 및 규제} & 정책 및 규제 준수 관리, 데이터 품질 관리, 보안 및 개인정보 보호 & 위험 관리, 편향 감지 및 완화, 공정성 및 윤리적 AI \\
\textbf{운영} & 감사 및 보고, 거버넌스 워크플로우, 데이터 리니지, 측정 & 투명성 및 설명 가능성, 규제 준수, AI 교육 및 훈련 \\
\textbf{통합} & 플랫폼 간 거버넌스, 데이터 공유, 데이터 통합 & AI 기술 평가, AI 스킬 평가, AI 거버넌스 통합 \\
\bottomrule
\end{tabular}
\end{adjustbox}
\captionof{table}{데이터 거버넌스 및 AI 거버넌스 역량 비교}
\label{tab:gov_capabilities}

\subsubsection*{데이터 관리 vs. 데이터 거버넌스}
이 두 개념은 밀접하게 관련되어 있지만, 초점이 다릅니다.
\begin{itemize}
    \item \textbf{데이터 관리 (Data Management):} 데이터의 운영적 측면, 즉 데이터를 어떻게 다룰 것인가(How)에 중점을 둡니다. 데이터 수명 주기(생성, 저장, 사용, 보관, 폐기) 전반에 걸친 기술적, 운영적 프로세스를 포함합니다.
    \item \textbf{데이터 거버넌스 (Data Governance):} 데이터 관리의 '누가(Who)', '무엇을(What)', '어떻게(How)'에 대한 정책, 절차, 표준을 수립하는 데 중점을 둡니다. 데이터 품질, 규제 준수, 보안 등 데이터 관리의 '규칙'을 정의합니다.
\end{itemize}

\subsubsection*{AI 관리 vs. AI 거버넌스}
데이터 관리와 유사하게 AI 관리와 AI 거버넌스도 구분됩니다.
\begin{itemize}
    \item \textbf{AI 관리 (AI Management):} AI 모델의 운영적 프로세스, 즉 모델을 어떻게 개발하고 배포하며 모니터링할 것인가(MLOps 수명 주기)에 중점을 둡니다.
    \item \textbf{AI 거버넌스 (AI Governance):} AI 모델의 '누가', '무엇을', '어떻게'에 대한 윤리적, 법적, 사회적 측면을 다룹니다. "이것을 해야 하는가?", "윤리적인가?", "어떤 파급 효과가 있을까?"와 같은 질문에 답하며, 단순히 프로세스를 운영하는 것을 넘어섭니다.
\end{itemize}

\subsubsection*{정보 거버넌스의 주요 목표 요약}
조직은 정보 거버넌스를 통해 다음과 같은 목표를 달성하고자 합니다.
\begin{itemize}
    \item \textbf{통합된 데이터 생태계 (Unified Data Ecosystem):} 파편화된 데이터를 통합하여 일관된 관리 환경을 구축합니다.
    \item \textbf{표준화된 데이터 관행 (Standardized Data Practices):} 조직 전체에 걸쳐 데이터 처리 및 사용에 대한 일관된 규칙을 적용합니다.
    \item \textbf{협업 증진 (Collaboration):} 데이터 사일로(silo)를 허물고 팀 간 데이터 공유 및 협업을 장려합니다.
    \item \textbf{데이터 리터러시 향상 (Increased Data Literacy):} 조직 구성원들이 데이터의 의미와 사용법을 쉽게 이해할 수 있도록 메타데이터, 비즈니스 친화적인 이름(시맨틱 레이어) 등을 제공합니다.
    \item \textbf{소유권 및 책임성 (Ownership and Accountability):} 모든 데이터 및 AI 자산에 대해 명확한 소유자와 책임자를 지정합니다.
\end{itemize}

\newpage
\subsection*{거버넌스 성숙도 모델 (Governance Maturity Models)}

\begin{keyconceptbox}
\textbf{한 줄 핵심 요약:} 거버넌스 성숙도 모델은 조직이 데이터 및 AI 거버넌스를 얼마나 잘 구축하고 실행하고 있는지 평가하고 개선하기 위한 단계별 프레임워크입니다.
\end{keyconceptbox}

\textbf{직관적 예시/비유:}
거버넌스 성숙도 모델은 아이가 성장하는 과정과 비슷합니다. 처음에는 아무것도 모르고(Initial/Aware), 점차 반복되는 행동을 인지하고(Managed/Reactive), 규칙을 배우고(Defined), 자신의 행동을 측정하며(Quantitatively Managed), 최종적으로는 스스로 최적화된 행동을 하는(Optimized/AI-driven) 단계로 나아갑니다. 조직도 이와 같이 거버넌스 역량을 점진적으로 발전시켜 나갑니다.

\textbf{기술적 정확한 설명:}
거버넌스 성숙도 모델은 조직이 데이터 및 AI 거버넌스 여정에서 현재 어느 위치에 있는지 평가하고, 개선이 필요한 영역을 식별하며, 다른 조직과 비교(벤치마킹)하여 체계적으로 역량을 향상시키는 데 도움을 줍니다.

\subsubsection*{거버넌스 성숙도 5단계}
\begin{adjustbox}{max width=\textwidth}
\begin{tabular}{p{0.15\textwidth}p{0.2\textwidth}p{0.6\textwidth}}
\toprule
\textbf{단계} & \textbf{특징} & \textbf{설명} \\
\midrule
\textbf{1. 초기/인지 (Initial/Aware)} & 개별적 노력, 비체계적 & 거버넌스에 대한 인식이 낮거나 부재한 단계. 스타트업처럼 당장 제품 출시나 고객 확보에 집중하여 거버넌스 표준이 거의 없는 상태입니다. \\
\textbf{2. 관리/반응적 (Managed/Reactive)} & 반복 작업 인지, 일부 규칙 수립 & 반복되는 데이터 관련 작업을 인지하고, 일부 표준화의 필요성을 느끼기 시작합니다. 문제가 발생했을 때 반응적으로 규칙을 수립하는 단계입니다. \\
\textbf{3. 정의됨 (Defined)} & 전사적 프레임워크, 이해관계자 동의 & 전사적인 데이터 거버넌스 프레임워크를 정의하고, 다양한 팀과 이해관계자의 동의를 얻어 규칙을 수립합니다. 많은 조직이 이 단계에서 어려움을 겪거나 멈춥니다. \\
\textbf{4. 정량적으로 관리됨 (Quantitatively Managed)} & 측정 및 개선, ROI 분석 & 정의된 거버넌스 프레임워크가 얼마나 효과적인지 측정하고, 그 결과를 바탕으로 개선합니다. 데이터 거버넌스의 ROI(투자 수익률)를 분석하고 위험을 완화합니다. \\
\textbf{5. 최적화/AI 기반 (Optimized/AI-driven)} & 자동화, 지속적 자가 개선 & 거버넌스 프로세스가 자동화되고, AI를 활용하여 지속적으로 스스로 개선되는 단계입니다. (예: Google Drive에서 민감 파일 공유 시 자동 경고 및 권한 조정) \\
\bottomrule
\end{tabular}
\end{adjustbox}
\captionof{table}{데이터 거버넌스 성숙도 5단계}
\label{tab:maturity_model}

\subsubsection*{AI 거버넌스 성숙도}
AI 거버넌스도 유사한 5단계 성숙도 모델을 따릅니다. 그러나 AI 거버넌스는 비교적 새로운 분야이므로, 대부분의 조직은 아직 3~4단계에 머물러 있으며, 5단계에 도달한 조직은 매우 드뭅니다. 이는 이 분야에 많은 성장 기회가 있음을 의미합니다.

\begin{cautionbox}
\textbf{성숙도 단계 건너뛰기 금지:}
조직은 성숙도 단계를 건너뛰어 빠르게 최적화 단계로 가려 할 수 있지만, 이는 실패할 가능성이 높습니다. 마치 아이가 성장하는 것처럼, 각 단계를 충분히 거쳐야 견고한 거버넌스 체계를 구축할 수 있습니다. 각 단계에서 충분히 성숙한 후 다음 단계로 나아가는 '폭포수 모델(waterfall model)' 접근 방식이 장기적으로 더 효과적입니다.
\end{cautionbox}

\newpage
\subsection*{위험도에 따른 거버넌스 수준}

\begin{keyconceptbox}
\textbf{한 줄 핵심 요약:} 데이터의 민감도와 잠재적 영향에 따라 거버넌스 통제의 엄격성 수준을 다르게 적용해야 합니다.
\end{keyconceptbox}

\textbf{직관적 예시/비유:}
집에 있는 물건들을 관리하는 것과 같습니다. 현금이나 귀금속(고위험 데이터)은 금고에 넣어두고 여러 잠금장치를 걸어두지만, 일반적인 생활용품(저위험 데이터)은 그냥 선반에 두거나 서랍에 넣어둡니다. 모든 물건을 똑같이 엄격하게 관리할 필요는 없으며, 위험도에 따라 관리 수준을 조절하는 것이 효율적입니다.

\textbf{기술적 정확한 설명:}
모든 데이터와 애플리케이션에 동일한 수준의 거버넌스 통제를 적용하는 것은 비효율적입니다. 데이터의 민감도, 잠재적 위험, 규제 요구사항 등을 고려하여 거버넌스 수준을 차등화해야 합니다.

\subsubsection*{위험도 결정 요인}
\begin{itemize}
    \item \textbf{산업별 및 규제 요구사항 (Industry-specific and Regulatory Requirements):} 금융, 의료 등 특정 산업은 엄격한 규제를 받습니다.
    \item \textbf{개인 및 조직에 미치는 영향 (Impact on Individuals and Organization):} 데이터 유출 시 개인의 피해(신용 정보 도용 등)와 조직의 평판 손상, 법적 책임 등을 고려합니다.
    \item \textbf{규모 (Scale):} 영향을 받는 개인의 수가 100명인지 100만 명인지에 따라 위험도가 달라집니다.
\end{itemize}

\subsubsection*{거버넌스 엄격성 스펙트럼}
\begin{figure}[h!]
    \centering
    \includegraphics[width=0.8\textwidth]{example-risk-spectrum.png} % Placeholder for a conceptual diagram
    \caption{위험도에 따른 거버넌스 엄격성 스펙트럼 (개념도)}
    \label{fig:risk_spectrum}
\end{figure}

\begin{itemize}
    \item \textbf{고위험 데이터 및 애플리케이션 (High-Risk Data and Applications):}
    \begin{itemize}
        \item \textbf{예시:} 민감한 개인 식별 정보(PII), 금융 데이터, 건강 정보.
        \item \textbf{거버넌스 수준:} 매우 엄격한 거버넌스(Strict Governance)가 필요합니다. 접근 권한 획득을 위해 여러 단계를 거쳐야 하며, 강력한 통제 메커니즘이 적용됩니다. 조직은 이러한 데이터가 침해될 경우 막대한 대가를 치러야 합니다.
    \end{itemize}
    \item \textbf{저위험 데이터 및 애플리케이션 (Low-Risk Data and Applications):}
    \begin{itemize}
        \item \textbf{예시:} 익명화되거나 토큰화된 데이터, PII가 없는 합성 데이터, 내부 머신러닝 실험 데이터.
        \item \textbf{거버넌스 수준:} 유연하고 신속한 거버넌스(Responsiveness)가 가능합니다. 내부적으로 누가 접근하든 큰 문제가 발생하지 않으므로, 접근 절차가 간소화될 수 있습니다.
    \end{itemize}
\end{itemize}
이 스펙트럼은 엄격한 통제와 비즈니스 민첩성 사이의 균형을 찾는 것을 목표로 합니다.

\newpage
\subsection*{일반적인 거버넌스 고려사항}

\begin{keyconceptbox}
\textbf{한 줄 핵심 요약:} 거버넌스는 단순히 도구를 구매하는 것으로 완성되지 않으며, 조직 전체의 사람, 프로세스, 기술이 통합되어 지속적으로 발전해야 합니다.
\end{keyconceptbox}

\textbf{직관적 예시/비유:}
건강한 식습관을 갖는 것과 같습니다. 비싼 유기농 식재료(최고의 거버넌스 도구)를 산다고 해서 저절로 건강해지는 것이 아닙니다. 식재료를 어떻게 요리하고(프로세스), 누가 요리하며(사람), 어떤 조리도구(기술)를 사용하는지, 그리고 꾸준히 실천하는지가 중요합니다.

\textbf{기술적 정확한 설명:}
\begin{itemize}
    \item \textbf{거버넌스는 구매할 수 없습니다 (You cannot buy data and governance):}
    단 하나의 도구를 구매하여 모든 거버넌스 문제를 해결할 수는 없습니다. 거버넌스는 \textbf{사람(People)}, \textbf{프로세스(Process)}, \textbf{기술(Technology)}의 세 가지 요소에 걸쳐 구현되어야 합니다.
    \item \textbf{모든 팀의 참여 (Depends on all teams):}
    거버넌스는 특정 부서만의 책임이 아니라, 데이터를 다루는 모든 팀의 '제2의 천성(second nature)'이 되어야 합니다. 이를 위해 교육과 리터러시 향상이 필수적이며, 상위 리더십의 강력한 지원과 지시가 필요합니다.
    \item \textbf{단순하게 시작 (Always start simple):}
    거버넌스 성숙도 모델의 최적화 단계부터 시작하려 하면 실패할 가능성이 높습니다. 작게 시작하여 점진적으로 개선하고 확장해나가야 합니다.
    \item \textbf{비즈니스 영역별 수준 선택 (Choosing the level):}
    조직 내 각 비즈니스 영역(예: 인사 부서 vs. 제품 개발 팀)의 특성과 데이터 민감도에 따라 거버넌스 수준과 여정이 다를 수 있습니다. 각 영역에 맞는 유연한 접근 방식이 필요합니다.
\end{itemize}

\newpage
\subsection*{거버넌스 역할 및 책임}

\begin{keyconceptbox}
\textbf{한 줄 핵심 요약:} 효과적인 거버넌스를 위해서는 정책을 정의하는 역할과 이를 실제 시스템에 구현하고 강제하는 역할이 명확하게 구분되어야 합니다.
\end{keyconceptbox}

\textbf{직관적 예시/비유:}
학교에서 규칙을 만드는 교장 선생님(정책 정의)과, 그 규칙을 학생들에게 설명하고 지키도록 지도하는 담임 선생님(정책 구현 및 강제)의 역할과 같습니다. 각자의 역할이 명확해야 학교가 질서 있게 운영될 수 있습니다.

\textbf{기술적 정확한 설명:}
조직 내에는 거버넌스 정책을 수립하고, 이를 실제 운영 환경에 적용하며, 준수 여부를 감독하는 다양한 역할이 존재합니다.

\subsubsection*{주요 거버넌스 역할 및 태스크}
\begin{adjustbox}{max width=\textwidth}
\begin{tabular}{p{0.2\textwidth}p{0.35\textwidth}p{0.35\textwidth}}
\toprule
\textbf{역할} & \textbf{데이터 거버넌스 태스크} & \textbf{AI 거버넌스 태스크} \\
\midrule
\textbf{정의 (Definition)} & 데이터 거버넌스 정책 개발 및 승인 & AI 거버넌스 정책 개발 및 승인 \\
\textbf{운영 (Operations)} & 데이터 거버넌스 정책 구현 및 강제 & AI 거버넌스 정책 구현 및 강제 \\
\midrule
\textbf{주요 역할} & & \\
\textbf{최고 거버넌스 책임자 (CGO)} & 전사적 거버넌스 전략 수립 및 감독 & AI 윤리 및 규제 준수 총괄 \\
\textbf{데이터 거버넌스 위원회} & 정책 및 표준 정의, 이해관계자 조정 & AI 모델 개발 및 배포 가이드라인 수립 \\
\textbf{데이터 스튜어드 (Data Steward)} & 데이터 정의, 품질 관리, 접근 권한 관리 & AI 모델 데이터셋 관리, 편향성 검토 \\
\textbf{데이터/AI 엔지니어} & 거버넌스 정책 기술적 구현, 시스템 구축 & AI 모델 배포, 모니터링, 보안 구현 \\
\textbf{법률 및 규제 준수 팀} & 법적 요구사항 해석 및 정책 반영 & AI 관련 법률 및 규제 준수 검토 \\
\bottomrule
\end{tabular}
\end{adjustbox}
\captionof{table}{주요 거버넌스 역할 및 태스크}
\label{tab:gov_roles}

\begin{tcolorbox}[cautionbox]
\textbf{AI 관련 역할의 부상:}
최근 AI 기술의 발전과 함께 '데이터 및 AI 스튜어드' 또는 'AI 스튜어드'와 같이 AI가 포함된 새로운 역할들이 등장하고 있습니다. 이는 AI 거버넌스의 중요성이 커지고 있음을 반영합니다.
\end{tcolorbox}

\subsubsection*{Ops의 중첩 (Overlapping Operations)}
조직 내에는 다양한 'Ops' 기능들이 존재하며, 이들은 거버넌스와 밀접하게 관련되어 있습니다.
\begin{itemize}
    \item \textbf{거버넌스 Ops (Governance Ops):} 거버넌스 정책 및 절차의 운영적 측면을 다룹니다.
    \item \textbf{데이터 Ops (Data Ops):} 데이터 수명 주기 관리 및 데이터 파이프라인의 운영적 효율성을 다룹니다.
    \item \textbf{MLOps (Machine Learning Ops):} 머신러닝 모델의 개발, 배포, 모니터링의 운영적 측면을 다룹니다.
\end{itemize}
이 세 가지 Ops는 서로 중첩되는 부분이 많습니다. 또한, 인사(People Ops), 영업(Sales Ops), 재무(FinOps)와 같은 사업부별 Ops도 각자의 데이터와 AI/ML 활용에 대한 거버넌스 과제를 안고 있습니다. 모든 Ops가 공통의 거버넌스 구조에 맞춰 표준화되면 효율성이 크게 증대되지만, 이를 위해서는 조직 전반의 동의와 가치 증명이 필요합니다.

\newpage
\section*{Databricks Lakehouse 아키텍처와 거버넌스}

\subsection*{Well-Architected Framework 확장}

\begin{keyconceptbox}
\textbf{한 줄 핵심 요약:} Databricks는 클라우드 공급업체의 기존 Well-Architected Framework에 '데이터 거버넌스'와 '상호운용성 및 유용성'이라는 두 가지 핵심 요소를 추가하여 Lakehouse 아키텍처를 위한 프레임워크를 확장했습니다.
\end{keyconceptbox}

\textbf{직관적 예시/비유:}
집을 짓는 데 필요한 기본적인 뼈대(클라우드 5대 원칙) 위에, 집의 안전과 규칙(데이터 거버넌스)을 담당하는 보안 시스템과, 다른 집들과의 연결성 및 사용 편의성(상호운용성 및 유용성)을 담당하는 현관 및 도로를 추가한 것과 같습니다. Databricks는 데이터 중심의 환경에서 이 두 가지가 특히 중요하다고 본 것입니다.

\textbf{기술적 정확한 설명:}
모든 주요 클라우드 공급업체(AWS, Azure, GCP)는 비용 최적화, 성능 효율성, 안정성, 보안, 운영 우수성이라는 5가지 핵심 원칙을 기반으로 하는 Well-Architected Framework를 제공합니다. Databricks는 이 5가지 원칙에 더해 Lakehouse 아키텍처의 특성을 반영한 두 가지 추가 원칙을 제시합니다.

\subsubsection*{Databricks Lakehouse Well-Architected Framework의 7가지 원칙}
\begin{enumerate}
    \item \textbf{운영 우수성 (Operational Excellence)}
    \item \textbf{보안 (Security)}
    \item \textbf{안정성 (Reliability)}
    \item \textbf{성능 효율성 (Performance Efficiency)}
    \item \textbf{비용 최적화 (Cost Optimization)}
    \item \textbf{데이터 거버넌스 (Data Governance):} Lakehouse 내의 모든 데이터 자산을 효과적으로 관리하고 보호하는 원칙입니다.
    \item \textbf{상호운용성 및 유용성 (Interoperability and Usability):} Databricks 외부의 데이터 소스(RDBMS, SaaS 등)와의 연동 및 통합 거버넌스를 포함하여, 데이터의 접근성과 활용성을 높이는 원칙입니다.
\end{enumerate}
Databricks는 이 두 가지 추가 원칙을 통해 고객들이 겪는 공통적인 거버넌스 및 통합 문제를 해결하고, Lakehouse 환경에서의 성숙도를 평가하는 프로그램을 제공합니다.

\newpage
\subsection*{Unity Catalog (UC) 개요}

\begin{keyconceptbox}
\textbf{한 줄 핵심 요약:} Unity Catalog는 Databricks Lakehouse 플랫폼의 핵심 거버넌스 계층으로, 모든 데이터 및 AI 자산에 대한 중앙 집중식 메타데이터 관리, 접근 제어, 감사, 리니지 기능을 제공합니다.
\end{keyconceptbox}

\textbf{직관적 예시/비유:}
Unity Catalog는 도서관의 중앙 관리 시스템과 같습니다. 어떤 책(데이터)이 어디에 있고, 누가 빌려갈 수 있으며, 언제 빌려갔고, 어떤 내용인지(메타데이터)를 모두 기록하고 관리합니다. 이 시스템을 통해 모든 책이 안전하게 관리되고, 필요한 사람이 쉽게 찾아볼 수 있으며, 도서관 전체의 규칙이 일관되게 적용됩니다.

\textbf{기술적 정확한 설명:}
Unity Catalog (UC)는 Databricks의 데이터 인텔리전스 플랫폼의 기반이 되는 거버넌스 솔루션입니다. Delta Lake, Iceberg, Parquet과 같은 정형 데이터뿐만 아니라 PDF, CSV 등 모든 유형의 비정형 데이터까지 클라우드 스토리지에 저장된 모든 자산을 UC를 통해 거버넌스할 수 있습니다. UC는 데이터에 접근하기 위한 필수적인 거버넌스 계층 역할을 하며, 이를 통해 그 위에 구축되는 모든 애플리케이션의 보안을 강화합니다.

\subsubsection*{해결하려는 문제}
UC는 조직이 겪는 다음과 같은 문제들을 해결하고자 합니다.
\begin{itemize}
    \item \textbf{파편화된 거버넌스 (Fragmented Governance):} 데이터가 여러 곳에 흩어져 있어 일관된 거버넌스 적용이 어렵습니다.
    \item \textbf{협업 부족 (Lack of Collaboration):} 데이터 공유가 어렵거나 허용되지 않아 팀 간 데이터 중복 및 프로젝트 중복이 발생합니다.
    \item \textbf{내장된 인텔리전스 부족 (Lack of Built-in Intelligence):} 데이터의 '진실의 원천(source of truth)'을 찾기 어렵고, 데이터 존재 여부조차 알기 힘듭니다.
\end{itemize}

\subsubsection*{Unity Catalog의 핵심 기능}
UC는 이러한 문제를 해결하기 위해 다음과 같은 기능을 제공합니다.
\begin{itemize}
    \item 접근 제어 (Access Control)
    \item 감사 (Auditing)
    \item 데이터 검색 (Discovery)
    \item 데이터 공유 (Data Sharing)
    \item 데이터 리니지 (Lineage)
    \item 데이터 품질 모니터링 (Quality Monitoring)
    \item 비용 통제 (Cost Controls)
\end{itemize}

\subsubsection*{UC Metastore의 진화}
\begin{adjustbox}{max width=\textwidth}
\begin{tabular}{p{0.45\textwidth}p{0.45\textwidth}}
\toprule
\textbf{기존 Metastore (UC 이전)} & \textbf{Unity Catalog Metastore} \\
\midrule
테이블 형식 데이터만 지원 (tabular data) & 모든 유형의 자산 지원 (테이블, 파일, 모델, 볼륨 등) \\
각 Databricks 워크스페이스에 사일로화됨 & 여러 Databricks 워크스페이스에 걸쳐 공통 \\
메타데이터가 분리되어 거버넌스 적용 어려움 & 중앙 집중식 메타데이터 계층으로 거버넌스 용이 \\
\bottomrule
\end{tabular}
\end{adjustbox}
\captionof{table}{기존 Metastore와 Unity Catalog Metastore 비교}
\label{tab:uc_metastore_comparison}

\subsubsection*{Unity Catalog의 기본 개념}
UC는 클라우드 환경의 기본 요소들을 추상화하여 거버넌스를 용이하게 합니다.
\begin{itemize}
    \item \textbf{자격 증명 (Credential):} AWS의 IAM 역할, Azure의 관리 ID, GCP의 서비스 계정과 같은 클라우드 인증 정보를 의미합니다.
    \item \textbf{외부 위치 (External Location):} S3, ADLS, GCS와 같은 클라우드 객체 스토리지 경로를 의미합니다.
    \item \textbf{볼륨 (Volume):} 외부 위치 내에 저장된 PDF 파일과 같은 비정형 데이터를 관리하는 데 사용됩니다.
    \item \textbf{외부 연결 (Foreign Connection):} Snowflake, BigQuery, Redshift 등 다른 RDBMS 또는 SaaS 플랫폼에 대한 연결을 의미합니다.
\end{itemize}

\subsubsection*{통합 거버넌스}
UC는 로컬 Databricks 테이블과 외부 연결을 통해 가져온 테이블에 동일한 거버넌스 원칙을 적용할 수 있도록 합니다.
\begin{lstlisting}[caption={Unity Catalog를 통한 통합 쿼리 예시}, label={lst:unified_query}, language=SQL]
-- Databricks 내 로컬 테이블 쿼리
SELECT * FROM main.paul.wine;

-- 외부 Snowflake 테이블 쿼리 (Lakehouse Federation으로 연결)
SELECT * FROM snowflake.schema.table;
\end{lstlisting}
위 예시처럼, 로컬 테이블과 외부 테이블 모두에 동일한 \texttt{GRANT} 문을 사용하여 접근 권한을 부여하고 관리할 수 있습니다. 이는 거버넌스 스토리를 하나의 플랫폼 아래로 통합하는 데 큰 이점을 제공합니다.

\subsubsection*{보안 계층}
UC는 클라우드 공급업체의 기본 보안(예: IAM 역할) 위에 추가적인 보안 계층을 제공합니다.
\begin{itemize}
    \item 클라우드 IAM 역할은 스토리지에 대한 광범위한 접근(예: \texttt{GET}, \texttt{PUT}, \texttt{LIST})을 가질 수 있습니다.
    \item 하지만 UC는 사용자에게 특정 데이터 객체에 대한 \texttt{SELECT} 권한만 부여하여, IAM 역할이 가진 \texttt{PUT} 권한이 있더라도 사용자가 \texttt{INSERT}를 할 수 없도록 제한할 수 있습니다.
\end{itemize}
이러한 2단계 보안 접근 방식은 사일로화된 IAM 역할의 필요성을 줄이고, 더 세분화된 접근 제어를 가능하게 합니다.

\subsubsection*{Unity Catalog 아키텍처}
\begin{figure}[h!]
    \centering
    \includegraphics[width=0.8\textwidth]{example-uc-arch.png} % Placeholder for a conceptual diagram
    \caption{Unity Catalog 아키텍처 (개념도)}
    \label{fig:uc_arch}
\end{figure}
\begin{itemize}
    \item 사용자는 컴퓨트 클러스터나 SQL 웨어하우스를 통해 Databricks에 접근합니다.
    \item UC 서비스는 Databricks 컨트롤 플레인에 위치하며, 접근 제어(Access Control)를 담당합니다.
    \item 사용자가 데이터를 쿼리하면, UC는 먼저 해당 사용자가 요청된 객체에 대한 접근 권한(예: \texttt{SELECT})을 가지고 있는지 확인합니다.
    \item 권한이 부여되면, UC는 클라우드 자격 증명(IAM 역할 등)에 데이터 요청을 위임하여 실제 클라우드 스토리지에서 데이터를 가져옵니다.
    \item 이 과정에서 사용자의 신원은 클라우드 스토리지 버킷에 직접 접근하지 않으며, UC가 연결을 중개합니다.
\item 이는 AWS의 '역할 가정(assuming a role)'과 유사한 개념으로, UC가 특정 IAM 역할을 가정하여 데이터에 접근하는 방식입니다.
\end{itemize}

\subsubsection*{데이터 저장 권장 사항}
UC에서는 클라우드 경로를 카탈로그 및 스키마 수준에서 컨테이너로 등록하는 것을 권장합니다.
\begin{itemize}
    \item 스키마 A의 데이터는 특정 경로에, 스키마 B의 데이터는 다른 경로에 저장하여 자연스러운 데이터 격리(data isolation)를 달성합니다.
    \item 각 스키마는 다른 IAM 역할로 거버넌스될 수 있어, 엄격한 데이터 격리 요구사항을 훨씬 간단하게 충족할 수 있습니다.
\end{itemize}

\newpage
\subsection*{속성 기반 접근 제어 (Attribute-Based Access Control, ABAC)}

\begin{keyconceptbox}
\textbf{한 줄 핵심 요약:} ABAC는 데이터 객체에 부여된 '태그'를 기반으로 접근 규칙을 정의하고, 이 규칙을 자동으로 적용하여 세분화된 데이터 보안을 구현하는 방식입니다.
\end{keyconceptbox}

\textbf{직관적 예시/비유:}
ABAC는 옷장 정리와 같습니다. 옷(데이터)에 '겨울옷', '여름옷', '정장', '캐주얼'과 같은 태그를 붙여둡니다. 그리고 "영희는 겨울옷만 입을 수 있다", "철수는 정장만 입을 수 있다"와 같은 규칙(정책)을 정합니다. 그러면 영희가 옷장을 열면 겨울옷만 보이고, 철수가 열면 정장만 보이는 식입니다. 데이터에 태그를 붙이고, 그 태그에 따라 누가 무엇을 볼 수 있는지 자동으로 결정하는 것입니다.

\textbf{기술적 정확한 설명:}
기존의 테이블 단위 또는 역할 기반 접근 제어(RBAC)만으로는 복잡한 조직의 거버넌스 요구사항을 충족하기 어렵습니다. ABAC는 데이터 요소에 '속성(Attribute)' 또는 '태그(Tag)'를 부여하고, 이 태그를 기반으로 데이터 접근 규칙을 정의하여 일관되고 자동화된 방식으로 적용합니다.

\subsubsection*{ABAC의 작동 방식}
\begin{enumerate}
    \item \textbf{컬럼 및 테이블 태그 지정 (Tag Columns and Tables):}
    데이터 내의 특정 컬럼이나 테이블에 'Govern Tags'라는 키-값 쌍을 부여합니다 (예: \texttt{PII: SSN}, \texttt{PII: Address}). 이 태그는 수동으로 부여할 수도 있고, Databricks의 데이터 분류(Data Classification) 기능을 통해 자동으로 부여될 수도 있습니다.
    \item \textbf{데이터 접근 규칙 정책 생성 (Create a Policy):}
    태그를 기반으로 데이터 접근 규칙을 정의하는 정책을 생성합니다. 예를 들어, "PII 태그가 있는 데이터는 '분석가 그룹'에게 마스킹된 버전으로 보여야 한다"와 같은 규칙을 설정합니다.
    \item \textbf{자동 적용 (Automatic Application):}
    정책이 생성되면, 해당 태그가 있는 데이터에 접근하는 사용자는 자동으로 정의된 규칙(예: 마스킹)이 적용된 데이터를 보게 됩니다.
\end{enumerate}

\subsubsection*{데이터 분류 (Data Classification)}
Databricks는 AI 기반 데이터 분류 기능을 제공하여 PII(개인 식별 정보)를 자동으로 감지하고 태그를 부여합니다.
\begin{itemize}
    \item \textbf{자동 감지:} 이메일 주소, 위치, 이름 등 PII 유형을 자동으로 식별합니다.
    \item \textbf{백그라운드 실행:} 데이터가 랜딩될 때마다 AI 모델이 증분 데이터에 대해 백그라운드에서 실행되어 자동으로 태그를 적용합니다.
    \item \textbf{시스템 태그:} 자동으로 부여된 태그는 '시스템 태그'로 관리됩니다.
    \item \textbf{이점:} 예를 들어, '이메일 주소'로 자동 분류된 모든 필드에 대해 '분석가 그룹'에게 마스킹 규칙을 적용하는 정책을 카탈로그 수준에서 한 번만 설정하면, 새로 추가되는 이메일 주소 필드에도 자동으로 규칙이 적용되어 거버넌스 부담을 크게 줄일 수 있습니다.
\end{itemize}

\subsubsection*{ABAC 정책 (ABAC Policy)}
ABAC 정책은 현재 컬럼 마스킹(Column Masking)과 행 필터링(Row Filtering)을 지원합니다.
\begin{itemize}
    \item \textbf{컬럼 마스킹:} 특정 컬럼의 데이터를 부분적으로 가리거나(예: SSN의 마지막 4자리만 표시) 완전히 가립니다.
    \item \textbf{행 필터링:} 특정 조건에 해당하는 행(row)을 데이터셋에서 제외하여 보여줍니다 (예: EU 고객 데이터 숨기기).
    \item \textbf{계층적 적용:} 카탈로그 및 스키마 수준에서 정책을 정의할 수 있으며, 상위 수준의 정책은 하위 수준으로 전파됩니다. 장기적으로는 계정 수준에서 모든 카탈로그에 적용되는 정책을 지원할 예정입니다.
\end{itemize}

\begin{examplebox}
\textbf{ABAC 정책 구현 예시:}
\begin{enumerate}
    \item \textbf{Govern Tag 생성:} \texttt{PII: SSN} 태그를 생성하고, 이 태그를 사용할 수 있는 권한을 부여합니다.
    \item \textbf{데이터 분류:} Databricks의 AI 기반 데이터 분류 기능이 자동으로 데이터셋 내의 SSN 컬럼을 감지하고 \texttt{PII: SSN} 태그를 부여합니다. (또는 수동으로 부여)
    \item \textbf{ABAC 정책 생성:}
    \begin{itemize}
        \item \textbf{정책 이름:} \texttt{MaskSSNForAnalysts}
        \item \textbf{적용 대상:} \texttt{analyst\_group}
        \item \textbf{범위:} 특정 카탈로그 내 모든 스키마
        \item \textbf{조건:} \texttt{PII: SSN} 태그가 있는 모든 컬럼
        \item \textbf{액션:} 컬럼 마스킹 함수 적용 (예: \texttt{substring(ssn, 6, 4)} 또는 \texttt{*****ssn})
    \end{itemize}
    \item \textbf{결과:} \texttt{analyst\_group}의 사용자가 해당 카탈로그 내의 어떤 테이블에서든 SSN 컬럼을 쿼리하면, 자동으로 마스킹된 SSN만 보게 됩니다.
\end{enumerate}
\end{examplebox}

\newpage
\subsection*{데이터 품질 모니터링 (Data Quality Monitoring)}

\begin{keyconceptbox}
\textbf{한 줄 핵심 요약:} 데이터 품질 모니터링은 Databricks Lakehouse에 유입되는 데이터의 신선도, 완전성, 정확성 등을 자동으로 감지하고 보고하여 데이터 신뢰도를 높이는 기능입니다.
\end{keyconceptbox}

\textbf{직관적 예시/비유:}
냉장고에 있는 음식(데이터)의 신선도를 자동으로 체크해주는 스마트 냉장고와 같습니다. "이 우유는 유통기한이 지났습니다!", "이 칸에 있던 채소가 갑자기 없어졌습니다!"와 같이 문제가 생기면 즉시 알려주어, 상한 음식을 먹거나 재료가 부족한 상황을 방지할 수 있습니다.

\textbf{기술적 정확한 설명:}
데이터 품질 모니터링은 거버넌스 관점에서 조직이 운영하는 데이터가 깨끗하고 신뢰할 수 있는지 확인하는 데 필수적입니다. Databricks는 이 기능을 통해 데이터 품질 규칙을 한 번만 설정하면, 백그라운드에서 자동으로 실행되어 데이터 품질 문제를 감지합니다.

\subsubsection*{주요 기능}
\begin{itemize}
    \item \textbf{이상 감지 (Anomaly Detection):} 데이터 패턴의 변화를 감지합니다. 예를 들어, 매일 1,000개의 행이 유입되던 테이블에 갑자기 0개의 행이 유입되면 이를 이상으로 감지하여 경고합니다.
    \item \textbf{데이터 프로파일링 (Data Profiling):} 데이터의 최소/최대값, 평균, 고유 값 수 등 통계 정보를 제공합니다.
    \item \textbf{AI 기반:} 백그라운드에서 AI 모델이 실행되어 데이터의 신선도(Freshness), 완전성(Completeness), 세분화(Segmentation) 등을 분석합니다.
    \item \textbf{결과 게시 및 알림:} 모든 모니터링 결과는 사용자가 선택한 테이블에 게시되므로, 이를 쿼리하여 대시보드를 구축하거나 알림 시스템을 연동할 수 있습니다.
    \item \textbf{간단한 설정:} 스키마 수준에서 한 번만 설정하면 해당 스키마 내의 모든 테이블에 적용됩니다. 향후에는 규칙을 사용자 정의할 수 있는 기능도 제공될 예정입니다.
\end{itemize}

\newpage
\subsection*{데이터 리니지 (Data Lineage)}

\begin{keyconceptbox}
\textbf{한 줄 핵심 요약:} 데이터 리니지는 데이터가 어디에서 시작하여 어떤 변환 과정을 거쳐 현재의 형태로 존재하는지, 그리고 어디로 흘러가는지를 시각적으로 추적할 수 있는 기능입니다.
\end{keyconceptbox}

\textbf{직관적 예시/비유:}
음식의 레시피와 재료 원산지를 추적하는 것과 같습니다. "이 요리(BI 보고서의 데이터)는 어떤 재료(원천 데이터)로 만들어졌고, 어떤 조리 과정(데이터 변환)을 거쳤으며, 어떤 식당(다른 시스템)으로 배달되었는가?"를 파악하는 것입니다. 문제가 생겼을 때 원인을 찾아내고 수정하는 데 필수적입니다.

\textbf{기술적 정확한 설명:}
데이터 리니지는 데이터 거버넌스에서 투명성과 책임성을 확보하는 데 매우 중요합니다. BI 보고서의 특정 데이터 요소가 어디에서 왔는지, 어떤 변환 로직을 거쳤는지 파악하여 데이터의 신뢰성을 검증하고 문제 발생 시 원인을 분석할 수 있습니다.

\subsubsection*{리니지 기능의 진화}
\begin{itemize}
    \item \textbf{초기 리니지:} Databricks 내에서 실행되는 모든 작업(job)에 대해 소스(source)와 타겟(target) 테이블 간의 리니지를 자동으로 추적했습니다. 하지만 Databricks 외부의 소스나 타겟에 대한 정보는 블랙박스였습니다.
    \item \textbf{사용자 정의 리니지 (Custom Lineage):} 최근에는 사용자가 Databricks 외부의 소스(예: Salesforce)나 타겟(예: PowerBI)을 직접 연결하여 리니지 그래프에 포함시킬 수 있는 기능이 추가되었습니다. 이를 통해 데이터의 전체 수명 주기에 걸친 엔드-투-엔드 리니지를 구축할 수 있습니다.
\end{itemize}

\subsubsection*{작동 방식}
\begin{itemize}
    \item Databricks에서 제출되는 모든 코드(SQL, Python 등)는 리니지 서비스에 의해 분석됩니다.
    \item 리니지 서비스는 컬럼과 테이블 간의 종속성을 파싱하여, 예를 들어 하나의 컬럼이 여러 컬럼으로 분할되거나(예: '이름' 컬럼이 '성', '이름', '중간 이름'으로 분할) 특정 문자열 함수가 적용되는 등의 변환 과정을 추적합니다.
    \item 리니지 정보는 API 및 시스템 테이블을 통해 외부 도구와도 통합될 수 있습니다.
\end{itemize}

\newpage
\subsection*{레이크하우스 페더레이션 (Lakehouse Federation)}

\begin{keyconceptbox}
\textbf{한 줄 핵심 요약:} 레이크하우스 페더레이션은 Databricks 외부의 다양한 데이터 소스(RDBMS, SaaS 등)를 Unity Catalog에 '연결'하여, 마치 Databricks 내의 데이터처럼 통합된 방식으로 쿼리하고 거버넌스할 수 있도록 하는 기능입니다.
\end{keyconceptbox}

\textbf{직관적 예시/비유:}
집에 있는 여러 개의 다른 언어로 된 책(외부 데이터베이스)들을 하나의 통합 번역기(레이크하우스 페더레이션)를 통해 읽는 것과 같습니다. 번역기를 통해 어떤 언어의 책이든 쉽게 읽을 수 있고, 누가 어떤 책을 읽을 수 있는지에 대한 규칙(거버넌스)도 번역기 시스템에서 일관되게 관리할 수 있습니다.

\textbf{기술적 정확한 설명:}
조직의 모든 데이터가 Databricks Lakehouse에만 존재하는 것은 아닙니다. Salesforce, Oracle, SQL Server, Snowflake 등 다양한 RDBMS 및 SaaS 플랫폼에 데이터가 분산되어 있습니다. 레이크하우스 페더레이션은 이러한 외부 소스에 대한 '연결(Connection)'을 Unity Catalog에 등록하여, Databricks 환경에서 이들 데이터를 직접 쿼리하고 거버넌스할 수 있도록 합니다.

\subsubsection*{주요 이점}
\begin{itemize}
    \item \textbf{통합 거버넌스:} 외부 소스의 테이블도 Unity Catalog에 등록된 델타 테이블과 동일한 \texttt{GRANT} 문을 사용하여 접근 권한을 관리할 수 있습니다. 이는 거버넌스 계층을 통합하여 관리 복잡성을 줄입니다.
    \item \textbf{쿼리 푸시다운 (Query Pushdown):} 대부분의 외부 소스에 대해 효율적인 쿼리 푸시다운을 지원합니다. 즉, Databricks가 데이터를 모두 가져와서 처리하는 대신, 쿼리의 일부 또는 전체를 외부 데이터베이스로 '밀어 넣어' 해당 데이터베이스에서 직접 처리하도록 합니다. 이는 네트워크 트래픽을 줄이고 쿼리 성능을 향상시킵니다.
    \item \textbf{다양한 소스 지원:} SQL Server, Salesforce, AWS Glue, Oracle, Redshift, PostgreSQL 등 다양한 RDBMS 및 클라우드 메타스토어 연결을 지원하며, 지원 소스는 지속적으로 확장됩니다.
\end{itemize}

\begin{cautionbox}
\textbf{Unity Catalog의 중요성:}
레이크하우스 페더레이션을 포함하여, 앞서 설명한 모든 Databricks의 거버넌스 관련 기능들은 Unity Catalog를 기반으로 합니다. Unity Catalog가 없으면 이러한 기능들을 활용할 수 없습니다. UC는 Databricks 플랫폼에서 거버넌스를 구현하기 위한 필수적인 기반 계층입니다.
\end{cautionbox}

\begin{tcolorbox}[summarybox]
\textbf{시장 동향:}
Snowflake의 Polaris/Horizon, Microsoft Fabric 등 모든 주요 데이터 플랫폼 공급업체는 자사 생태계 내의 모든 자산에 대한 거버넌스 계층을 구축하고 있습니다. 궁극적으로는 다른 플랫폼의 자산까지 가장 효과적으로 거버넌스할 수 있는 솔루션이 시장에서 우위를 점할 것으로 예상됩니다.
\end{tcolorbox}

\newpage
\subsection*{델타 쉐어링 (Delta Sharing)}

\begin{keyconceptbox}
\textbf{한 줄 핵심 요약:} 델타 쉐어링은 Databricks Lakehouse의 데이터, 노트북, AI 모델 등 다양한 자산을 클라우드나 플랫폼에 관계없이 안전하고 개방적인 방식으로 공유할 수 있도록 하는 기술입니다.
\end{keyconceptbox}

\textbf{직관적 예시/비유:}
파일을 이메일로 보내거나 USB에 담아주는 대신, 클라우드 기반의 안전한 공유 링크를 통해 친구에게 파일을 공유하는 것과 같습니다. 친구는 어떤 기기나 운영체제를 사용하든 그 링크를 통해 파일을 열어볼 수 있으며, 파일 소유자는 누가 언제 파일을 열어봤는지 관리할 수 있습니다.

\textbf{기술적 정확한 설명:}
델타 쉐어링은 Databricks가 주도하는 오픈 소스 프로토콜이자 관리형 서비스입니다. Unity Catalog와 긴밀하게 통합되어 UC에 의해 거버넌스되는 모든 자산을 동일한 거버넌스 규칙을 적용하여 공유할 수 있도록 합니다.

\subsubsection*{주요 특징}
\begin{itemize}
    \item \textbf{다양한 자산 공유:} 데이터셋(Delta 테이블), 노트북, AI 모델, 볼륨 등 Databricks 내의 다양한 자산을 공유할 수 있습니다.
    \item \textbf{Unity Catalog 통합:} UC의 거버넌스 규칙(접근 제어, 감사 등)이 델타 쉐어링에도 그대로 적용됩니다.
    \item \textbf{크로스 클라우드/플랫폼 공유:}
    \begin{itemize}
        \item \textbf{클라우드 간 공유:} AWS, Azure, GCP 등 다른 클라우드 환경에 있는 사용자에게 데이터를 공유할 수 있습니다.
        \item \textbf{플랫폼 간 공유:} Databricks 사용자가 아니더라도 Tableau, PowerBI, Pandas 등 다양한 도구에서 공유된 데이터를 소비할 수 있습니다.
    \end{itemize}
    \item \textbf{레이크하우스 페더레이션과의 시너지:} 레이크하우스 페더레이션을 통해 Unity Catalog에 등록된 Snowflake 테이블과 같은 외부 데이터도 델타 쉐어링을 통해 다른 클라우드나 플랫폼의 사용자에게 공유할 수 있습니다.
    \item \textbf{개방형 프로토콜:} 델타 쉐어링은 개방형 프로토콜이므로, Databricks 고객이 아니어도 데이터를 소비할 수 있습니다.
\end{itemize}

\begin{tcolorbox}[summarybox]
\textbf{시장 경쟁:}
Snowflake, Redshift 등 다른 데이터 플랫폼들도 자체적인 데이터 공유 기술을 가지고 있습니다. 누가 더 편리하고 다양한 기능을 제공하며, 다른 플랫폼과의 상호운용성을 잘 지원하는지가 시장 경쟁의 핵심 요소입니다.
\end{tcolorbox}

\newpage
\subsection*{상호운용성 (Interoperability)}

\begin{keyconceptbox}
\textbf{한 줄 핵심 요약:} Databricks는 Delta Lake뿐만 아니라 Iceberg와 같은 다른 주요 데이터 형식과의 상호운용성을 강화하여, 고객이 원하는 데이터 형식을 자유롭게 사용하고 공유할 수 있도록 지원합니다.
\end{keyconceptbox}

\textbf{직관적 예시/비유:}
다양한 종류의 스마트폰(데이터 플랫폼)이 서로 다른 운영체제(데이터 형식)를 사용하더라도, 공통의 충전기(상호운용성)를 사용하여 충전하고 데이터를 주고받을 수 있도록 하는 것과 같습니다. 어떤 기기를 사용하든 데이터가 막히지 않고 흐를 수 있도록 하는 것입니다.

\textbf{기술적 정확한 설명:}
데이터 레이크하우스 생태계에서는 Delta Lake와 Iceberg가 두 가지 주요 데이터 형식으로 자리 잡고 있습니다. Databricks는 자사의 Delta Lake 기술을 넘어 Iceberg 테이블도 지원하고 공유할 수 있도록 함으로써, 고객이 특정 데이터 형식에 얽매이지 않고 유연하게 데이터를 관리하고 공유할 수 있는 환경을 제공합니다. 이는 데이터 생태계를 더욱 강력하고 개방적으로 만듭니다.

\newpage
\section*{Databricks UI를 통한 거버넌스 실습 (데모 요약)}

강의에서는 Databricks UI를 통해 Unity Catalog의 거버넌스 기능을 시연했습니다. 다음은 주요 시연 내용과 그 의미를 요약한 것입니다.

\subsection*{카탈로그 탐색기 및 ABAC 정책}

\begin{enumerate}
    \item \textbf{카탈로그 탐색:} Databricks 워크스페이스의 '카탈로그 탐색기'에서 특정 카탈로그를 선택합니다. 엔터프라이즈 환경에서는 수많은 카탈로그와 스키마가 존재합니다.
    \item \textbf{정책 탭 (Policies Tab):} 카탈로그 또는 스키마 수준에서 '정책(Policies)' 탭을 확인할 수 있습니다.
    \item \textbf{정책 적용 확인:}
    \begin{itemize}
        \item 카탈로그 수준에서 정의된 정책(예: \texttt{hide EU customers}, \texttt{mask SSN})은 하위 스키마로 상속됩니다.
        \item 스키마 수준에서도 해당 스키마에만 적용되는 정책을 추가로 생성할 수 있습니다.
        \item 장기적으로는 계정 수준에서 모든 카탈로그에 적용되는 정책을 지원할 예정입니다.
    \end{itemize}
    \item \textbf{새 정책 생성 (New Policy):}
    \begin{itemize}
        \item \textbf{이름 지정:} 정책의 이름을 지정합니다.
        \item \textbf{적용 대상 (Apply to):} 이 정책의 적용을 받을 사용자 또는 그룹을 지정합니다 (예: \texttt{all account users}).
        \item \textbf{예외 (Except):} '브레이크 글라스(break glass)' 접근 방식으로, 특정 사용자나 그룹은 정책의 적용을 받지 않고 모든 데이터를 볼 수 있도록 예외를 설정할 수 있습니다. 이는 테스트나 감사 시 유용합니다.
        \item \textbf{범위 (Scope):} 정책을 적용할 카탈로그, 스키마, 테이블을 지정합니다. 특정 태그가 있는 테이블에만 적용하도록 설정할 수 있습니다.
        \item \textbf{조건 (Custom Conditions):} 태그 값 등을 조합하여 복잡한 조건을 설정할 수 있습니다.
        \item \textbf{액션 (Mask/Hide):}
        \begin{itemize}
            \item \textbf{마스킹 (Mask):} 특정 컬럼을 마스킹하는 함수를 정의합니다.
            \item \textbf{숨기기 (Hide):} 행 필터링 함수를 정의하여 특정 행을 숨깁니다.
        \end{itemize}
        \item \textbf{Databricks Assistant 활용:} 마스킹/숨기기 함수 작성 시 Databricks Assistant의 도움을 받을 수 있습니다.
    \end{itemize}
    \item \textbf{마스킹 정책 예시 (\texttt{mask SSN}):}
    \begin{itemize}
        \item \textbf{적용 대상:} \texttt{all account users} (예외 없음).
        \item \textbf{조건:} \texttt{MGMPI} 태그(사용자 정의 태그)가 있는 컬럼.
        \item \textbf{함수:} SSN을 마스킹하는 함수(예: \texttt{*****6789})를 정의하고 테스트합니다.
    \end{itemize}
    \item \textbf{Govern Tags 생성 및 적용:}
    \begin{itemize}
        \item 'Govern Tags' 페이지에서 \texttt{PII}와 같은 태그 키와 \texttt{SSN}, \texttt{Address}와 같은 허용 값을 정의합니다.
        \item 특정 사용자에게 이 태그를 '할당(assign)'할 수 있는 권한을 부여합니다.
        \item \textbf{데모:} \texttt{clickstream} 테이블의 \texttt{current\_page\_title} 컬럼에 \texttt{MGMPI: SSN} 태그를 수동으로 적용했습니다.
    \end{itemize}
    \item \textbf{태그 적용에 따른 마스킹 변화 시연:}
    \begin{itemize}
        \item \texttt{current\_page\_title} 컬럼에 \texttt{MGMPI: SSN} 태그를 적용한 후 테이블을 쿼리하자, 해당 컬럼의 데이터가 마스킹되어 표시되었습니다. (실제 SSN이 아니었지만, 태그와 정책의 연동을 보여주기 위함)
        \item 태그를 삭제한 후 다시 쿼리하자, 데이터가 마스킹되지 않고 원본 그대로 표시되었습니다.
        \item 이는 ABAC가 태그, 정책, 규칙을 기반으로 계층적으로 작동하며, 효과적으로 구현될 경우 매우 강력한 거버넌스 도구가 될 수 있음을 보여줍니다.
    \end{itemize}
\end{enumerate}

\subsection*{데이터 거버넌스 페이지}

Databricks는 조직의 거버넌스 현황을 한눈에 파악할 수 있는 '데이터 거버넌스(Data Governance)' 페이지를 제공합니다.

\begin{enumerate}
    \item \textbf{개요 (Overview):}
    \begin{itemize}
        \item 전체 데이터 자산에 대한 통계(예: 활성 테이블 수, 사용된 테이블 수).
        \item 인기 있는 테이블 태그, 컬럼 태그 등을 시각적으로 보여줍니다.
        \item 주석이 없는 테이블 목록을 제공하여 데이터 리터러시 개선을 위한 팀의 노력을 유도할 수 있습니다.
        \item 테이블에 대한 \texttt{SELECT} 권한, \texttt{MODIFY} 권한 등을 가진 사용자 현황을 다양한 방식으로 분석할 수 있습니다.
    \end{itemize}
    \item \textbf{데이터 분류 (Data Classification):}
    \begin{itemize}
        \item 자동으로 스캔된 테이블 목록과 감지된 PII 요소(예: 이메일 주소, 위치)를 보여줍니다.
        \item 특정 PII 유형(예: 이메일 주소)을 선택하면, 해당 PII를 포함하는 모든 테이블 목록을 확인할 수 있습니다.
        \item '자동 태그(auto-tag)' 기능을 통해 감지된 PII에 자동으로 태그를 적용할 수 있습니다.
        \item '시스템 태그(system tag)'는 Databricks가 자동으로 분류하여 부여한 태그로, 스패너 아이콘으로 표시됩니다.
    \end{itemize}
    \item \textbf{데이터 품질 모니터링 (Data Quality Monitoring):}
    \begin{itemize}
        \item 모니터링 중인 스키마 및 테이블의 '건강 상태(health status)'를 보여줍니다 (건강함/건강하지 않음).
        \item '건강하지 않음'으로 표시된 테이블의 경우, 어떤 문제가 있는지(예: \texttt{commit freshness is stale} - 커밋 신선도 저하) 상세한 정보를 제공합니다.
        \item 데이터 프로파일링(Data Profiling)을 구성하여 시계열 프로파일링, 스냅샷 프로파일링 등 다양한 통계 정보를 수집할 수 있습니다.
        \item 모든 모니터링 결과는 사용자가 지정한 테이블에 게시되므로, 이를 기반으로 맞춤형 대시보드를 구축하거나 알림 시스템을 연동할 수 있습니다.
    \end{itemize}
\end{enumerate}

\subsection*{레이크하우스 페더레이션 연결}

\begin{enumerate}
    \item \textbf{연결 탭 (Connections Tab):} Databricks 워크스페이스의 '연결(Connections)' 탭에서 다양한 외부 데이터 소스에 대한 연결 목록을 확인할 수 있습니다.
    \item \textbf{지원되는 소스:} SQL Server, Salesforce, AWS Glue, Oracle, Amazon Redshift, PostgreSQL 등 다양한 RDBMS 및 클라우드 메타스토어에 대한 연결이 생성되어 있음을 시연했습니다.
    \item \textbf{통합 거버넌스:} 이러한 외부 연결을 통해 가져온 모든 데이터 객체는 Unity Catalog에 등록되며, Databricks 내의 델타 테이블과 동일한 거버넌스 원칙(예: \texttt{GRANT} 문)을 적용하여 관리할 수 있습니다. 이는 파편화된 데이터 환경에서 거버넌스 계층을 통합하는 강력한 방법입니다.
\end{enumerate}

\newpage
\section*{FAQ}

초심자가 거버넌스 개념을 학습하면서 혼동할 수 있는 질문들을 Q\&A 형식으로 정리했습니다.

\begin{tcolorbox}[mybox, title={Q: 데이터 유출은 거버넌스 다이어그램에 어떻게 포함되나요?}]
A: 데이터 유출 사고는 직접적으로 다이어그램의 한 요소로 표시되지는 않지만, '법률 및 규제 준수'와 '조직에 미치는 영향'이라는 측면에서 거버넌스 정책 수립의 강력한 동기가 됩니다. 유출 시 발생하는 금전적 벌금, 평판 손상, 고객 신뢰 상실 등의 위험을 피하기 위해 조직은 거버넌스를 강화합니다.
\end{tcolorbox}

\begin{tcolorbox}[mybox, title={Q: AI 거버넌스 없이 데이터 거버넌스만 할 수 있나요?}]
A: 네, 할 수 있습니다. 사실 AI 거버넌스는 데이터 거버넌스가 선행되어야 효과적입니다. AI 모델은 데이터에 기반하므로, 데이터의 품질, 보안, 적법성이 확보되지 않으면 AI 모델의 윤리적 사용이나 책임성을 논하기 어렵습니다. 데이터 거버넌스는 AI 거버넌스를 위한 필수적인 기반입니다.
\end{tcolorbox}

\begin{tcolorbox}[mybox, title={Q: 거버넌스 성숙도 단계를 건너뛸 수 있나요?}]
A: 이론적으로는 가능하지만, 실제로는 매우 어렵고 실패할 가능성이 높습니다. 각 단계는 이전 단계에서 얻은 경험과 기반 위에 구축되므로, 단계를 건너뛰면 나중에 예상치 못한 문제에 직면할 수 있습니다. 장기적인 관점에서는 각 단계를 충분히 거쳐 견고한 기반을 다지는 것이 더 효과적입니다.
\end{tcolorbox}

\begin{tcolorbox}[mybox, title={Q: 거버넌스 도구 하나로 모든 것을 해결할 수 있나요?}]
A: 불가능합니다. 시장에 존재하는 어떤 단일 도구도 모든 거버넌스 요구사항을 충족할 수는 없습니다. 거버넌스는 '사람, 프로세스, 기술'의 통합적인 노력이며, 여러 최적의 도구들을 선택하고 이들을 서로 통합하여 사용하는 전략이 필요합니다. 도구의 난립을 피하면서도 핵심 기능을 충족하는 2~3개의 도구를 통합하는 것이 일반적입니다.
\end{tcolorbox}

\begin{tcolorbox}[mybox, title={Q: 케이스 스터디를 위한 조직은 가상이어도 되나요?}]
A: 강의 Q\&A에 따르면, 케이스 스터디를 위한 조직은 실제 존재하는 회사여야 합니다. 가상의 조직보다는 실제 회사를 바탕으로 그 회사의 구체적인 도전 과제와 특성을 분석하여 벤더 및 솔루션을 선택하는 것이 요구됩니다.
\end{tcolorbox}

\newpage
\section*{빠르게 훑어보기 (1페이지 요약)}

\begin{tcolorbox}[summarybox]
\textbf{데이터 및 AI 거버넌스 핵심 요약}

\begin{itemize}
    \item \textbf{정보 거버넌스:} 모든 정보 자산(데이터, AI, 물리적 문서)을 관리하는 최상위 틀. 법규 준수, 가치 극대화, 위험 완화가 목표.
    \item \textbf{데이터 vs. AI 거버넌스:}
    \begin{itemize}
        \item \textbf{데이터 거버넌스:} 데이터 수집, 저장, 보안 등 전통적 관리.
        \item \textbf{AI 거버넌스:} AI 시스템 개발, 배포, 윤리, 편향, 책임 관리 (데이터 거버넌스 위에 구축).
    \end{itemize}
    \item \textbf{거버넌스 성숙도 모델 (5단계):} 초기 $\rightarrow$ 관리 $\rightarrow$ 정의 $\rightarrow$ 정량적 관리 $\rightarrow$ 최적화. 점진적 발전이 중요하며, AI 거버넌스는 아직 초기 단계.
    \item \textbf{위험도 기반 거버넌스:} 민감도(PII, 금융)에 따라 거버넌스 엄격성 조절 (고위험 = 엄격, 저위험 = 유연).
    \item \textbf{Databricks Lakehouse Framework:} 클라우드 5대 원칙에 \textbf{데이터 거버넌스} 및 \textbf{상호운용성/유용성} 추가.
    \item \textbf{Unity Catalog (UC):} Databricks의 중앙 거버넌스 계층.
    \begin{itemize}
        \item 모든 데이터/AI 자산(테이블, 파일, 모델) 메타데이터 관리.
        \item 클라우드 기본 보안 위에 추가 보안 계층 제공.
        \item \textbf{ABAC (속성 기반 접근 제어):} 데이터 태그(예: PII) 기반으로 접근 규칙 정의 및 자동 적용.
        \item \textbf{데이터 분류:} AI로 PII 자동 감지 및 태그 부여.
        \item \textbf{데이터 품질 모니터링:} 데이터 신선도, 완전성 이상 감지 및 보고.
        \item \textbf{데이터 리니지:} 데이터 출처 및 변환 과정 추적 (외부 소스/타겟 연결 가능).
        \item \textbf{레이크하우스 페더레이션:} 외부 DB/SaaS를 UC에 연결하여 통합 쿼리 및 거버넌스.
        \item \textbf{델타 쉐어링:} UC 자산(데이터, 노트북, 모델)을 클라우드/플랫폼에 관계없이 안전하게 공유.
    \end{itemize}
    \item \textbf{핵심 원칙:} 거버넌스는 도구 구매만으로 완성되지 않으며, '사람, 프로세스, 기술'의 통합적 노력과 지속적인 개선이 필요.
\end{itemize}
\end{tcolorbox}

\newpage
\section*{학습 로드맵}

이 강의 내용을 효과적으로 학습하고 실제 역량으로 발전시키기 위한 로드맵입니다.

\begin{enumerate}
    \item \textbf{기초 개념 이해 (강의 내용 숙지):}
    정보 거버넌스, 데이터 거버넌스, AI 거버넌스의 정의, 목적, 구성 요소(정책, 절차, 표준, 통제)를 명확히 이해합니다. 각 개념이 왜 필요한지, 어떤 문제를 해결하는지 파악하는 데 집중합니다.
    \item \textbf{성숙도 모델 및 위험 관리:}
    거버넌스 성숙도 5단계와 위험도에 따른 거버넌스 수준 조절 원칙을 이해하여, 조직의 현재 위치를 파악하고 개선 방향을 설정하는 데 필요한 사고방식을 익힙니다.
    \item \textbf{Databricks Unity Catalog 핵심 기능:}
    Unity Catalog가 제공하는 ABAC, 데이터 분류, 품질 모니터링, 리니지, 페더레이션, 델타 쉐어링 등 주요 기능들이 각각 어떤 거버넌스 과제를 해결하는지 구체적으로 학습합니다. 각 기능의 작동 방식과 이점을 파악합니다.
    \item \textbf{실습 및 적용 (Databricks 환경):}
    가능하다면 Databricks Community Edition 또는 실제 워크스페이스에서 Unity Catalog를 설정하고, ABAC 정책, Govern Tags, 데이터 분류 결과 확인, 데이터 품질 모니터링 설정 등을 직접 실습해봅니다. 이론을 실제 환경에 적용하는 경험을 쌓습니다.
    \item \textbf{심화 학습 및 최신 동향:}
    AI 거버넌스의 윤리적 측면, 최신 규제 동향(예: AI Act), 그리고 다른 클라우드 벤더(Snowflake, Microsoft Fabric)의 거버넌스 솔루션에 대해 추가적으로 학습하여 폭넓은 시야를 확보합니다.
\end{enumerate}

\end{document}